%% ma: L12-B14-0143
%% lop: 12
%% chuong: 5
%% bai: 14
%% dang: 12.14.20
%% loai: TLN
%% muc: VD
%% dap_an: "7"
%% nguon: "Ảnh giáo viên gửi 25/11/2026 (tài liệu ôn tập Phương trình mặt phẳng), Phần III Câu 3 (ba điểm A(1;0;-2), B(-3;1;1), C(5;5;-5))"
%% kien_thuc: "Bài 14 (phương trình mặt phẳng đi qua ba điểm không thẳng hàng)"
%% trang_thai: chinh-thuc
\begin{ex}
Trong không gian $Oxyz$, biết mặt phẳng $(P)$ đi qua ba điểm $A(1;0;-2)$, $B(-3;1;1)$, $C(5;5;-5)$ có phương trình dạng $ax+by+cz+5=0$, tính $a+b+c$.
\shortans{7}
\loigiai{$\overrightarrow{AB}=(-4;1;3)$, $\overrightarrow{AC}=(4;5;-3)$, $\left[\overrightarrow{AB},\overrightarrow{AC}\right]=(-18;0;-24)=-6(3;0;4)$. Mặt phẳng $(P)$: $3(x-1)+4(z+2)=0\Leftrightarrow3x+4z+5=0$. Vậy $a=3$, $b=0$, $c=4$ và $a+b+c=7$.}
\end{ex}
